\documentclass{article}

%% page layout
\textheight 9in 
\textwidth 6.5in
\topmargin -0.5in
\oddsidemargin 0in
\evensidemargin 0in
%\hoffset=-0.5in
\voffset=-0.25in

%% packages
\usepackage{scrtime} % for \thistime (this package MUST be listed first!)
\usepackage{amsmath} % essential for cases environment
\usepackage{amsthm} % for theorems and proofs
\usepackage{amsfonts} % mathbb
\usepackage{graphics,graphicx}
\usepackage{multirow} % fancy tables
\usepackage{wasysym} % circle symbols (including half-filled circles)
\usepackage{enumerate} % fancier enumeration (e.g., a,b,c, ...)
%\usepackage{xcolor}
\usepackage{color}
\newcommand{\de}[1]{{\color{red}{\bfseries DE:} #1}}

%% for solutions to multiple choice questions:
\newcommand{\correct}{{\color{blue}\fbox{\color{red}\checkmark} }}

%% macros
\newcommand{\reals}{\mathbb{R}}
\newcommand{\term}[1]{{\bfseries\slshape #1}}
\newcommand{\Ker}{{\text{Ker}\,}}
\newcommand{\Range}{{\text{Range}\,}}
\newcommand{\diag}{{\text{diag}}}
\newcommand{\alg}{{\text{alg}}}
\newcommand{\geom}{{\text{geom}}}
\newcommand{\norm}[1]{\left\|#1\right\|}
\newcommand{\abs}[1]{\left|#1\right|}
\newcommand{\R}{{\cal R}}
\newcommand{\G}{{\cal G}}
\newcommand{\eps}{\varepsilon}
\newcommand{\B}{\cal B}
\newcommand{\Tinf}{T_\textrm{inf}}
\newcommand{\Shat}{{\hat{S}}}
\newcommand{\Ihat}{{\hat{I}}}
\newcommand{\ie}{\emph{i.e., }}
\newcommand{\eg}{\emph{e.g., }}
\newcommand{\Rlogo}{\protect\includegraphics[height=2ex,keepaspectratio]{images/Rlogo.pdf}\xspace}
\newcommand{\XPPAUT}{\texttt{XPPAUT}\xspace}
\newcommand{\etal}{\textit{et al}.\xspace}
\newcommand\emphblue[1]{\emph{\color{blue}#1}}

%%%%%%%%%%%%%%%%%%%%%%%%%%%%
%% from feverpreamble.tex %%
%%%%%%%%%%%%%%%%%%%%%%%%%%%%

\usepackage{amssymb,latexsym,amsmath,setspace}
%%\usepackage[colorlinks,linkcolor=blue]{hyperref}
\usepackage[colorlinks=true,allcolors=blue]{hyperref}
%%\usepackage[colorlinks]{hyperref}
\usepackage{xspace}
\usepackage{graphics,graphicx}
\usepackage{subfigure}
\usepackage{lineno}
\usepackage{fancyhdr}
\usepackage[english]{babel}  %% for texi2dvi ~ bug
\usepackage[normalem]{ulem}
\usepackage{tikz} % http://www.texample.net/tikz/examples/tikzdevice-demo/
  % N.B. version 0.6.3 of tikzDevice from Rforge is required!!
%% improve figure caption typsetting:  (see ~/tex/caption.pdf for manual)
\usepackage[footnotesize,bf]{caption}
\usepackage{placeins} % \FloatBarrier

% citation macros
\newcommand{\citen}[1]{\cite{#1}}

% comment macros
\usepackage{color}
\newcommand{\comment}[3]{\textcolor{#1}{\textbf{[#2: }\textit{#3}\textbf{]}}}
\newcommand{\david}[1]{\comment{blue}{DE}{#1}}
\newcommand{\paul}[1]{\comment{cyan}{PA}{#1}}
\newcommand{\ben}[1]{\comment{magenta}{BB}{#1}}
\newcommand{\needref}{{\textcolor{red}{[NEED REF]}}}
\newcommand{\TBD}{{\textcolor{red}{{\bf TBD}}}}
% remove comments
%\renewcommand{\comment}[3]{\relax}
% remove section headings
%\def\subsubsection*#1{\relax\nobreak}

% other macros
\newcommand{\avg}[1]{{\left\langle#1\right\rangle}}
\newcommand{\var}[1]{\textrm{var}\left(#1\right)}
\newcommand{\sem}[1]{\textrm{sem}\left(#1\right)}
\newcommand{\natinf}{{\mathcal I}}
\newcommand{\find}{f_{\textrm{i}}}
\newcommand{\fpop}{f_{\textrm{p}}}
\newcommand{\logit}{\textrm{logit}}
\newcommand{\sign}{\textrm{sign}}
\newcommand{\logistic}{\textrm{logistic}}
\def\AJE{{\it American Journal of Epidemiology\/}}
\newcommand{\code}[1]{{\tt #1}}
\newcommand{\magcode}[1]{{\tt\color{magenta}#1}}
\newcommand{\redcode}[1]{{\tt\color{red}#1}}
\newcommand{\blackcode}[1]{{\tt\color{black}#1}}

%%%%%%%%%%%%%%%%%%%
%% JOURNAL NAMES %%
%%%%%%%%%%%%%%%%%%%
\def\PNAS{PNAS}
\def\JAMA{JAMA}
\def\BMB{{\it Bulletin of Mathematical Biology\/}}

% references
\newcommand{\eref}[1]{Equation~\eqref{E:#1}}
\newcommand{\fref}[1]{Figure~\ref{F:#1}}
\newcommand{\tref}[1]{Table~\ref{T:#1}}
\newcommand{\sref}[1]{\S\ref{S:#1}}
% other macros
\newcommand{{\Reff}}{{\mathcal{R}}_{\rm eff}}
\newcommand{{\Sinit}}{S_{\rm init}}
\newcommand{\supp}{Supplementary Information}
\newcommand{\StoppedHere}{\bigskip\bigskip{\textcolor{red}{\hrule\centerline{\bfseries STOPPED HERE}\hrule}}\bigskip\bigskip}
\newcommand{\colvec}[2]{\begin{pmatrix}#1\\#2\end{pmatrix}}
\newcommand{\diagmat}[3]{\begin{pmatrix}#1&0&0\\0&#2&0\\0&0&#3\end{pmatrix}}

\newcommand{\solution}[1]{{\hfill\break\vspace{-0.5\baselineskip}\break\color{blue}\emph{Solution: }#1}}
\newcommand{\tr}{\text{tr}}

\newcommand{\thickredline}{\bigskip{\color{red}\hrule height 5pt}\bigskip}

%% for assignment 3:
\newtheorem{theorem}{Theorem}
\newtheorem{remark}{Remark}
\newcommand{\openset}{{\mathcal O}}
\newcommand{\C}{{\mathcal C}}

% Journal Names
% -------------
% \def\MNRAS{{\it Mon.\ Not.\ R.\ astr.\ Soc.\/}}
\def\MNRAS{{\it Monthly Notices of the Royal Astronomical Society\/}}
% \def\ApJ{{\it Astrophys.~J.\/}}
\def\ApJ{{\it The Astrophysical Journal\/}}
% \def\Interface{{\it J.\ R.\ Soc.\ Lond.\/ \rm Interface}}
\def\Interface{{\it Journal of the Royal Society of London, \rm Interface}}
% \def\ProcA{{\it Proc.\ R.\ Soc.\ Lond.\/ \rm A}}
\def\ProcA{{\it Proceedings of the Royal Society of London, Series\/ \rm A}}
% \def\ProcB{{\it Proc.\ R.\ Soc.\ Lond.\/ \rm B}}
\def\ProcB{{\it Proceedings of the Royal Society of London, Series\/ \rm B}}
% \def\TransB{{\it Phil.\ Trans.\ R.\ Soc.\ Lond.\/ \rm B}}
\def\TransB{{\it Philosophical Transactions of the Royal Society of London, Series\/ \rm B}}
\def\TREE{{\it Trends in Ecology and Evolution\/}}
% \def\JTB{{\it J.\ theor.\ Biol.\/}}
\def\JTB{{\it Journal of Theoretical Biology\/}}
% \def\BE{{\it Behav.\ Ecol.\/}}
\def\BE{{\it Behavioural Ecology\/}}
% \def\BES{{\it Behav.\ Ecol.\ Sociobiol.\/}}
\def\BES{{\it Behavioural Ecology and Sociobiology\/}}
% \def\BJLS{\it Biol.\ J.\ Linn.\ Soc.\/}}
\def\BJLS{{\it Biological Journal of the Linnean Society\/}}
\def\Nature{{\it Nature\/}}
\def\Science{{\it Science\/}}
\def\Lancet{{\it The Lancet\/}}
\def\LancetID{{\it Lancet Infectious Diseases\/}}
% \def\PNAS{{\it Proc.\ Natl.\ Acad.\ Sci.\ USA}}
\def\PNAS{{\it PNAS -- Proceedings of the National Academy of Sciences of the
U.S.A.}}
\def\PLoSMed{{\it PLoS Medicine\/}}
\def\PLoSCB{{\it PLoS Computational Biology\/}}
\def\EID{{\it Emerging Infectious Diseases\/}}
\def\BMB{{\it Bulletin of Mathematical Biology\/}}
\def\AJE{{\it American Journal of Epidemiology\/}}
\def\TPB{{\it Theoretical Population Biology\/}}
\def\JRSS{{\it Journal of the Royal Statistical Society\/}}
\def\JRSSB{{\it Journal of the Royal Statistical Society Series B\/}}
\def\IRV{{\it Influenza and Other Respiratory Viruses\/}}
\def\JAMA{{\it JAMA -- Journal of the American Medical Association\/}}
\def\JGLR{{\it Journal of Great Lakes Research\/}}
\def\NEJM{{\it New England Journal of Medicine\/}}
\def\JMB{{\it Journal of Mathematical Biology\/}}
\def\Interface{{\it Journal of the Royal Society Interface\/}}
\def\THEE{{\it Theoretical Ecology\/}}
\def\Annals{{\it Annals of Internal Medicine\/}}
\def\PRL{{\it Physical Review Letters\/}}
\def\PRX{{\it Physical Review X\/}}
\def\NAMS{{\it Notices of the American Mathematical Society\/}}

%% underline with smash through:
\newcommand*{\undersmash}[1]{\underline{\smash{#1}}}

%% referring to TeX macros
\newcommand\ttbackslash{{\tt\char`\\}}
\newcommand{\macro}[1]{{\tt\ttbackslash#1}}


\begin{document}

\section{$\R_0$ for small pox}

\subsection*{(a)}

Denote the susceptible state by $S$, the incubation state by $E$, and the recovered state by $R$. Furthermore, denote the four infectious stages by $I_1, I_2, I_3$, and $I_4$, respectively, and let the each of the stages represent the four different levels of infectiousness exhibited by infected individuals. For example, the $I_1$ stage corresponds to the prodrom stage, in which infectiousness is rare, while $I_4$ stage represents both the pustules \& scabs stage and the resolving scabs stage, as both of these stages share a low level of infectiousness. Note that all states are repsented by proportions.

Let $\sigma$ represent the per capita rate at which an infected individual develops symptoms and let $\gamma_i$ represent the per capita rate at which an infected individual in stage $i$ progresses to the next stage. In other words, $\gamma_i$ is considered to be the removal rate from stage $I_i$, such that the individual progresses to an infectiousness level exhibited in $I_{i+1}$. Moreover, note that $\gamma_4$ corresponds to the per capita recovery rate. 
Finally, let $\beta_i$ represent the infectiousness (per contact transmission rate) of an infected individual in stage $i$ and $\mu$ represent the per capita natural birth/death rate (they are assumed to be equal). Then, we can write an ODE model for this system as follows:
$$
\begin{aligned}
\frac{dS}{dt} &= \mu (1-S) - (\beta_1 I_1 + \beta_2 I_2 + \beta_3 I_3 + \beta_4 I_4) S\\
\frac{dE}{dt} &= (\beta_1 I_1 + \beta_2 I_2 + \beta_3 I_3 + \beta_4 I_4) S - (\sigma + \mu) E \\
\frac{dI_1}{dt} &= \sigma E - (\gamma_1 + \mu) I_1\\
\frac{dI_2}{dt} &= \gamma_1 I_1 - (\gamma_2 + \mu) I_1\\
\frac{dI_3}{dt} &= \gamma_2 I_2 - (\gamma_3 + \mu) I_1\\
\frac{dI_4}{dt} &= \gamma_3 I_3 - (\gamma_4 + \mu) I_1\\
\frac{dR}{dt} &= \gamma_1 I_4 - \mu R\\
\end{aligned}
$$

\subsection*{(b)}

In order for an infected individual to infect a susceptible individual, it must survive the incubation period and become infectious. 
Then, we can think of $\R_0$ as a sum of the basic reproductive numbers from each stage, such that the final number of secondary cases caused by a typical infective individual is merely equal to the number of infections caused by an infective at each stage. For example, the contribution of infection from the first stage would be equivalent to $\R_0$ of an SEIR model. We denote this value as $\R_{0_1}$:
$$
\beta_1 \times \frac{\sigma}{\sigma + \mu} \times \frac{1}{\gamma_1 + \mu}
$$
where $\beta_1/(\gamma_1 + \mu)$ represent the average number of infections that occur in stage 1 and $\sigma/(\sigma+\mu)$ is the probability that an infected individual does not die before the incubation period is over.
In order for infection to occur in stage $i>1$, an infected individual must not die from natural mortality before reaching stage $i$. Then, the contribution of infection during the second stage would be defined as $\R_{0_2}$, such that
$$
\beta_2 \times \frac{\sigma}{\sigma + \mu} \times \frac{\gamma_1}{\gamma_1 + \mu} \times \frac{1}{\gamma_2 + \mu}.
$$
Note that we now have $\gamma_1/(\gamma_1+\mu)$ to account for probability the of progressing to stage 2 without dying from natural causes while in stage 1.
Likewise, we can do a similar computation for all other stages. Summing each of our values for $\R_{0_i}$, we ultimately obtain
$$
\begin{aligned}
\R_0 &= \beta_1 \times \frac{\sigma}{\sigma + \mu} \times \frac{1}{\gamma_1 + \mu}\\
&+ \beta_2 \times \frac{\sigma}{\sigma + \mu} \times \frac{\gamma_1}{\gamma_1 + \mu} \times \frac{1}{\gamma_2 + \mu}\\
&+ \beta_3 \times \frac{\sigma}{\sigma + \mu} \times \frac{\gamma_1}{\gamma_1 + \mu} \times \frac{\gamma_2}{\gamma_2 + \mu} \times \frac{1}{\gamma_3 + \mu}\\
&+ \beta_4 \times \frac{\sigma}{\sigma + \mu} \times \frac{\gamma_1}{\gamma_1 + \mu} \times \frac{\gamma_2}{\gamma_2 + \mu} \times \frac{\gamma_3}{\gamma_3 + \mu} \times \frac{1}{\gamma_4 + \mu}.
\end{aligned}
$$

\subsection*{(c)}

For this particular system, we have
$$
\mathcal{F} = \begin{pmatrix}
(\beta_1 I_1 + \beta_2 I_2 + \beta_3 I_3 + \beta_4 I_4) S\\
0\\
0\\
0\\
0
\end{pmatrix}, \mathcal{V} = \begin{pmatrix}
(\sigma + \mu) E \\
- \sigma E + (\gamma_1 + \mu) I_1\\
- \gamma_1 I_1 + (\gamma_2 + \mu) I_1\\
- \gamma_2 I_2 + (\gamma_3 + \mu) I_1\\
- \gamma_3 I_3 + (\gamma_4 + \mu) I_1\\
\end{pmatrix}
$$
Linearizing at the disease free equilibrium, we have
$$
F = \begin{pmatrix}
0 & \beta_1 & \beta_2 & \beta_3 & \beta_4\\
0 & 0 & 0 & 0 & 0\\
0 & 0 & 0 & 0 & 0\\
0 & 0 & 0 & 0 & 0\\
0 & 0 & 0 & 0 & 0\\
\end{pmatrix},
V = \begin{pmatrix}
(\sigma + \mu) & 0 & 0 & 0 & 0\\
- \sigma & \gamma_1 + \mu & 0 & 0 & 0\\
0 & - \gamma_1 & \gamma_2 + \mu & 0 & 0\\
0 & 0 & - \gamma_2 & \gamma_3 + \mu & 0\\
0 & 0 & 0 & - \gamma_3 & \gamma_4 + \mu\\
\end{pmatrix}
$$
Inverting this matrix by hand, we observe:
$$
\begin{aligned}
&\left(\begin{array}{ccccc|ccccc}
(\sigma + \mu) & 0 & 0 & 0 & 0 & 1 & 0 & 0 & 0 & 0\\
- \sigma & \gamma_1 + \mu & 0 & 0 & 0 & 0 & 1 & 0 & 0 & 0\\
0 & - \gamma_1 & \gamma_2 + \mu & 0 & 0 & 0 & 0 & 1 & 0 & 0\\
0 & 0 & - \gamma_2 & \gamma_3 + \mu & 0 & 0 & 0 & 0 & 1 & 0\\
0 & 0 & 0 & - \gamma_3 & \gamma_4 + \mu & 0 & 0 & 0 & 0 & 1\\
\end{array}\right)\\
\implies
&\left(\begin{array}{ccccc|ccccc}
1 & 0 & 0 & 0 & 0 & \frac{1}{\sigma + \mu} & 0 & 0 & 0 & 0\\
- \sigma & \gamma_1 + \mu & 0 & 0 & 0 & 0 & 1 & 0 & 0 & 0\\
0 & - \gamma_1 & \gamma_2 + \mu & 0 & 0 & 0 & 0 & 1 & 0 & 0\\
0 & 0 & - \gamma_2 & \gamma_3 + \mu & 0 & 0 & 0 & 0 & 1 & 0\\
0 & 0 & 0 & - \gamma_3 & \gamma_4 + \mu & 0 & 0 & 0 & 0 & 1\\
\end{array}\right)\\
\implies
&\left(\begin{array}{ccccc|ccccc}
1 & 0 & 0 & 0 & 0 & \frac{1}{\sigma + \mu} & 0 & 0 & 0 & 0\\
0 & 1 & 0 & 0 & 0 & \frac{\sigma}{(\sigma + \mu)(\gamma_1 + \mu)} & \frac{1}{\gamma_1 + \mu} & 0 & 0 & 0\\
0 & - \gamma_1 & \gamma_2 + \mu & 0 & 0 & 0 & 0 & 1 & 0 & 0\\
0 & 0 & - \gamma_2 & \gamma_3 + \mu & 0 & 0 & 0 & 0 & 1 & 0\\
0 & 0 & 0 & - \gamma_3 & \gamma_4 + \mu & 0 & 0 & 0 & 0 & 1\\
\end{array}\right)\\
\implies
&\left(\begin{array}{ccccc|ccccc}
1 & 0 & 0 & 0 & 0 & \frac{1}{\sigma + \mu} & 0 & 0 & 0 & 0\\
0 & 1 & 0 & 0 & 0 & \frac{\sigma}{(\sigma + \mu)(\gamma_1 + \mu)} & \frac{1}{\gamma_1 + \mu} & 0 & 0 & 0\\
0 & 0 & 1 & 0 & 0 & \frac{\sigma \gamma_1}{(\sigma + \mu)(\gamma_1 + \mu)(\gamma_2 + \mu)} & \frac{\gamma_1}{(\gamma_1 + \mu)(\gamma_2 + \mu)} & \frac{1}{\gamma_2 + \mu} & 0 & 0\\
0 & 0 & - \gamma_2 & \gamma_3 + \mu & 0 & 0 & 0 & 0 & 1 & 0\\
0 & 0 & 0 & - \gamma_3 & \gamma_4 + \mu & 0 & 0 & 0 & 0 & 1\\
\end{array}\right)\\
\implies
&\left(\begin{array}{ccccc|ccccc}
1 & 0 & 0 & 0 & 0 & \frac{1}{\sigma + \mu} & 0 & 0 & 0 & 0\\
0 & 1 & 0 & 0 & 0 & \frac{\sigma}{(\sigma + \mu)(\gamma_1 + \mu)} & \frac{1}{\gamma_1 + \mu} & 0 & 0 & 0\\
0 & 0 & 1 & 0 & 0 & \frac{\sigma \gamma_1}{(\sigma + \mu)(\gamma_1 + \mu)(\gamma_2 + \mu)} & \frac{\gamma_1}{(\gamma_1 + \mu)(\gamma_2 + \mu)} & \frac{1}{\gamma_2 + \mu} & 0 & 0\\
0 & 0 & 0 & 1 & 0 & \frac{\sigma \gamma_1 \gamma_2}{(\sigma + \mu)(\gamma_1 + \mu)(\gamma_2 + \mu)(\gamma_3 + \mu)} & \frac{\gamma_1 \gamma_2}{(\gamma_1 + \mu)(\gamma_2 + \mu)(\gamma_3 + \mu)} & \frac{\gamma_2}{(\gamma_2 + \mu)(\gamma_3 + \mu)} & \frac{1}{\gamma_3 + \mu} & 0\\
0 & 0 & 0 & - \gamma_3 & \gamma_4 + \mu & 0 & 0 & 0 & 0 & 1\\
\end{array}\right)\\
\implies
&\tiny\left(\begin{array}{ccccc|ccccc}
1 & 0 & 0 & 0 & 0 & \frac{1}{\sigma + \mu} & 0 & 0 & 0 & 0\\
0 & 1 & 0 & 0 & 0 & \frac{\sigma}{(\sigma + \mu)(\gamma_1 + \mu)} & \frac{1}{\gamma_1 + \mu} & 0 & 0 & 0\\
0 & 0 & 1 & 0 & 0 & \frac{\sigma \gamma_1}{(\sigma + \mu)(\gamma_1 + \mu)(\gamma_2 + \mu)} & \frac{\gamma_1}{(\gamma_1 + \mu)(\gamma_2 + \mu)} & \frac{1}{\gamma_2 + \mu} & 0 & 0\\
0 & 0 & 0 & 1 & 0 & \frac{\sigma \gamma_1 \gamma_2}{(\sigma + \mu)(\gamma_1 + \mu)(\gamma_2 + \mu)(\gamma_3 + \mu)} & \frac{\gamma_1 \gamma_2}{(\gamma_1 + \mu)(\gamma_2 + \mu)(\gamma_3 + \mu)} & \frac{\gamma_2}{(\gamma_2 + \mu)(\gamma_3 + \mu)} & \frac{1}{\gamma_3 + \mu} & 0\\
0 & 0 & 0 & 0 & 1    & \frac{\sigma \gamma_1 \gamma_2 \gamma_3}{(\sigma + \mu)(\gamma_1 + \mu)(\gamma_2 + \mu)(\gamma_3 + \mu)(\gamma_4  +\mu)} & \frac{\gamma_1 \gamma_2 \gamma_3}{(\gamma_1 + \mu)(\gamma_2 + \mu)(\gamma_3 + \mu)(\gamma_4  +\mu)} & \frac{\gamma_2 \gamma_3}{(\gamma_2 + \mu)(\gamma_3 + \mu)(\gamma_4  +\mu)} & \frac{\gamma_3}{(\gamma_3 + \mu)(\gamma_4  +\mu)} & \frac{1}{\gamma_4  +\mu}\\
\end{array}\right)\\
\end{aligned}
$$
Then,
$$
V^{-1} = \begin{pmatrix}
\frac{1}{\sigma + \mu} & 0 & 0 & 0 & 0\\
\frac{\sigma}{(\sigma + \mu)(\gamma_1 + \mu)} & \frac{1}{\gamma_1 + \mu} & 0 & 0 & 0\\
\frac{\sigma \gamma_1}{(\sigma + \mu)(\gamma_1 + \mu)(\gamma_2 + \mu)} & \frac{\gamma_1}{(\gamma_1 + \mu)(\gamma_2 + \mu)} & \frac{1}{\gamma_2 + \mu} & 0 & 0\\
\frac{\sigma \gamma_1 \gamma_2}{(\sigma + \mu)(\gamma_1 + \mu)(\gamma_2 + \mu)(\gamma_3 + \mu)} & \frac{\gamma_1 \gamma_2}{(\gamma_1 + \mu)(\gamma_2 + \mu)(\gamma_3 + \mu)} & \frac{\gamma_2}{(\gamma_2 + \mu)(\gamma_3 + \mu)} & \frac{1}{\gamma_3 + \mu} & 0\\
\frac{\sigma \gamma_1 \gamma_2 \gamma_3}{(\sigma + \mu)(\gamma_1 + \mu)(\gamma_2 + \mu)(\gamma_3 + \mu)(\gamma_4  +\mu)} & \frac{\gamma_1 \gamma_2 \gamma_3}{(\gamma_1 + \mu)(\gamma_2 + \mu)(\gamma_3 + \mu)(\gamma_4  +\mu)} & \frac{\gamma_2 \gamma_3}{(\gamma_2 + \mu)(\gamma_3 + \mu)(\gamma_4  +\mu)} & \frac{\gamma_3}{(\gamma_3 + \mu)(\gamma_4  +\mu)} & \frac{1}{\gamma_4  +\mu}\\
\end{pmatrix}
$$
It is clear that matrix $FV^{-1}$ consists of 0 entries except its first row.
Hence, its eigenvalues are on its diagonal, four of which are zero. The only non-zero entry on the diagonal is the first column entry of the first row, which is equal to the previously derived $\R_0$ value:
$$
\begin{aligned}
&\beta_1 \times \frac{\sigma}{\sigma + \mu} \times \frac{1}{\gamma_1 + \mu}\\
&+ \beta_2 \times \frac{\sigma}{\sigma + \mu} \times \frac{\gamma_1}{\gamma_1 + \mu} \times \frac{1}{\gamma_2 + \mu}\\
&+ \beta_3 \times \frac{\sigma}{\sigma + \mu} \times \frac{\gamma_1}{\gamma_1 + \mu} \times \frac{\gamma_2}{\gamma_2 + \mu} \times \frac{1}{\gamma_3 + \mu}\\
&+ \beta_4 \times \frac{\sigma}{\sigma + \mu} \times \frac{\gamma_1}{\gamma_1 + \mu} \times \frac{\gamma_2}{\gamma_2 + \mu} \times \frac{\gamma_3}{\gamma_3 + \mu} \times \frac{1}{\gamma_4 + \mu}\\
\end{aligned}
$$
Therefore, the derivation of $\R_0$ using the next generation method is consistent with the derivation of $\R_0$ from a biological argument.

\subsection*{(d)}

Note that for unaltered small pox, the time scale of disease is much shorter than average life span of a person. Then, we can approximate $\R_0$ by assuming that $\mu \approx 0$. The expression for $\R_0$ thus becomes
$$
\R_0 \approx \frac{\beta_1}{\gamma_1} + \frac{\beta_2}{\gamma_2} + \frac{\beta_3}{\gamma_3} + \frac{\beta_4}{\gamma_4}
$$
The alteration causes the early rash stage to be twice as long and so $\gamma_2^{-1}$ changes from 4 days to 8 days.
It is evident that increasing $\gamma_2^{-1}$ will lead to increase in $\R_0$.

Provided that infectiousness during the early rash stage is extreme, we can assume that at least half of the infection occurs during this stage. In the worst case scenario, all infections occur during the early rash stage. These assumptions are reasonable given that disease-induced death can occur in later stages, and so there would be little contribution to infection.
Based on these assumptions, we have that
$$
2.5 < \frac{\beta_2}{\gamma_{2,\tiny{\textrm{original}}}} < 5.
$$
Since altering doubles $\gamma_2^{-1}$, we get
$$
5 < \frac{\beta_2}{\gamma_{2,\tiny{\textrm{altered}}}} < 10.
$$
This translates to
$$
7.5 < \R_{0, \tiny{\textrm{altered}}} <10.
$$

\subsection*{(e)}

We highly recommend CDC to start preparing for the smallpox terror immediately. Given that smallpox has been altered to have twice as long early rash stage, we expect $\R_0$ to increase by at least a factor of 1.5. In the worse case scenario, $\R_0$ of the altered smallpox with be twice as large as that of the unaltered smallpox, corresponding to the following range: $7.5 < \R_0 < 10$.
This means that the final size of an epidemic will be almost 100\% if the attack is successful.




\end{document}
